\chapter{Minimal Witnesses in Higher Dimensions}
\label{ch:higher-d-witnesses}

Parts I through III extended the reconstruction toward standard quantum
mechanics, chapter by chapter, until Chapter~\ref{ch:final-audit} closed that
arc. This part does something different: it no longer moves toward recovering
more of the textbook formalism, but stress-tests what has already been built,
asking how far the same admissibility machinery pushes on its own terms.
The governing questions become "how little is actually enough?" and "where
does this specific reconstruction genuinely run out?" — a research programme
extending the argument, not an appendix collecting loose ends.

Chapter~\ref{ch:minimal-witnesses} left open what $m_\beta(\R,\C;d)$ is for
$d\ge3$, noting that Chapter~\ref{ch:unistochastic}'s $n\ge3$ incompleteness
theorem gives a single-matrix separation whose relation to the graph-level
minimal witness was unresolved. This chapter resolves it — and the answer
inverts the $d=2$ picture rather than extending it.

\begin{theorem}[Collapse of the minimal witness for $d\ge3$]
\label{thm:higher-d-collapse}
\index{Collapse of the minimal witness for d>=3 (theorem)}
For every $d \ge 3$, $m_\beta(\R,\C;d) = 0$: a single edge, carrying no
cycle at all, already witnesses $\mathcal R_\R(G,d) \subsetneq \mathcal
R_\C(G,d)$. This is conditional on nothing beyond
Definitions~\ref{def:realizability-space}--\ref{def:minimal-witness} as
already stated; no new assumption is introduced.
\end{theorem}

\begin{proof}
Let $P_2$ be the two-vertex, one-edge graph ($\beta_1(P_2) = 1-2+1 = 0$).
By Theorem~\ref{thm:n3incompleteness} (the $n\ge3$ incompleteness theorem),
there is a unistochastic $\Gamma \in \mathrm{US}(3)\setminus\mathrm{OS}(3)$
— the modulus-square of the $3\times3$ DFT matrix. Label the single edge of
$P_2$ with $\Gamma$. Since $\Gamma \notin \mathrm{OS}(3)$, no real frame
assignment reproduces it: this edge-labeling is not in $\mathcal R_\R(P_2,3)$
(indeed it fails even edgewise real-admissibility, and by
Proposition~\ref{prop:tree-no-constraint} — global realizability on a tree
is equivalent to edgewise admissibility — failing the latter is exactly
failing the former). Since $\Gamma \in \mathrm{US}(3)$ by construction, it
\emph{is} realizable over $\C$: the labeling is in $\mathcal R_\C(P_2,3)$.
Hence $P_2$ is a witness with $\beta_1(P_2)=0$, giving $m_\beta(\R,\C;3)
\le 0$, and $\beta_1(G)\ge0$ for every graph, so $m_\beta(\R,\C;3)=0$
exactly. The same $\Gamma$, embedded as one edge of any tree on $d=3$
contexts for $d\ge3$ (padding the unused dimensions trivially), gives the
identical argument for every $d\ge3$.
\end{proof}

\begin{remark}[Why this inverts, rather than extends, the $d=2$ result]
Chapter~\ref{ch:minimal-witnesses}'s proof that $m_\beta(\R,\C;2)\ge1$ used
Proposition~\ref{prop:tree-no-constraint} together with a fact specific to
$d=2$: $\mathrm{US}(2)=\mathrm{OS}(2)$ (Theorem~\ref{thm:n2degeneracy}), so
every edge individually admissible over $\C$ at $d=2$ is automatically
admissible over $\R$ too, and no tree can separate the fields. That
equality is exactly what Theorem~\ref{thm:n3incompleteness} shows fails at
$d=3$. The $d=2$ case is the exceptional one: cycles are needed there
\emph{because} single edges cannot yet distinguish the fields. From $d=3$
onward, single edges already can, and the graph-theoretic machinery of
Chapters~\ref{ch:context-graphs}--\ref{ch:minimal-witnesses} — cycle rank,
spanning trees, the triangle-gluing theorem — is not needed to separate
$\R$ from $\C$ at all. It remains exactly what is needed for the sharper,
narrower question that motivated it: whether \emph{individually admissible}
data (data that passes every edge test) can still fail to glue globally.
That question is nontrivial at every $d\ge2$, and is not addressed by
Theorem~\ref{thm:higher-d-collapse}.
\end{remark}

\begin{remark}[What Theorem~\ref{thm:higher-d-collapse} does not say]
\label{rem:minimum-not-classification}
Theorem~\ref{thm:higher-d-collapse} establishes a statement about a
\emph{minimum}:
\begin{equation}
\min\{\text{witness size}\} = \begin{cases} 3, & d=2, \\ 1, & d \ge 3.
\end{cases}
\end{equation}
It does not establish that graph-theoretic gluing obstructions disappear
for $d\ge3$, and should not be read that way. The question left open is
whether there exist \emph{edgewise-admissible} labelings on cyclic graphs
— data that individually passes every single-edge test, so that no
Theorem~\ref{thm:higher-d-collapse}-style witness is available — that
nevertheless fail to admit a global realization, generalizing the $d=2$
MUB triangle phenomenon of Chapter~\ref{ch:real-mub-obstruction} to
$d\ge3$. This is a logically independent question from the one this
chapter answers: the cheaper single-edge obstruction found here does not
preclude a separate, more expensive cycle-level obstruction also existing
alongside it. If future work shows no such edgewise-admissible-but-
globally-inadmissible examples exist at $d\ge3$, Theorem
\ref{thm:higher-d-collapse} would then collapse the entire obstruction
theory to single-edge tests. If such examples are found instead, the
theory acquires two independent layers — local (single-edge) and global
(cycle-gluing) obstructions — the way many areas of mathematics
distinguish satisfying every local constraint from admitting a global
solution. Neither possibility is resolved here.
\end{remark}

\begin{ledgerentry}[End of Chapter~\ref{ch:higher-d-witnesses}]
\textbf{Established:} $m_\beta(\R,\C;d)=0$ for all $d\ge3$
(Theorem~\ref{thm:higher-d-collapse}), resolving Chapter
\ref{ch:minimal-witnesses}'s first open question; the reason is identified
precisely as the failure of $\mathrm{US}(d)=\mathrm{OS}(d)$ beyond $d=2$,
not a new phenomenon requiring new machinery.\\
\textbf{Not established:} anything about $m_\beta(\C,\Hset;d)$ (quaternionic
separation, Chapter~\ref{ch:minimal-witnesses}'s second open question,
addressed only partially in Chapter~\ref{ch:quaternionic}); the
density-threshold question, still untouched (Chapter
\ref{ch:minimal-witnesses}'s fourth open question resolves neither here nor
in Chapter~\ref{ch:forbidden-subgraph}, which settles the forbidden-subgraph
question at $d=2$ but not the density question at any $d$). The cycle-based
machinery of Chapters~\ref{ch:context-graphs}--\ref{ch:minimal-witnesses}
remains the right tool for the gluing question it was built for, and this
chapter's result should not be read as making that machinery unnecessary —
only as showing it answers a different, narrower question than "does $\R$
ever fail," which single edges already settle from $d=3$ on (see
Remark~\ref{rem:minimum-not-classification} above for exactly what remains
open about that narrower question).
\end{ledgerentry}

\begin{remark}[Two distinct hazards, not one]
This chapter and Chapter~\ref{ch:final-audit}'s discussion of recent
real-valued quantum mechanics literature share a family resemblance worth
naming precisely rather than merging. Both are cases where something that
looked like a fundamental obstruction turned out to be contingent — but on
different kinds of contingency, and the difference matters for what a
reader should check.

The 2025--2026 papers concern \emph{representation-dependence}: the
Kronecker product and Hoffreumon and Woods' $\otimes_r$ both describe the
identical physical composition, and the apparent obstruction to real
quantum theory belonged to a choice of notation, not to the physics. The
diagnostic is recoordinatization: does the obstruction survive a change of
representation for the \emph{same} object?

Theorem~\ref{thm:higher-d-collapse} concerns a different hazard: a
\emph{low-dimensional coincidence}. $\mathrm{US}(2)=\mathrm{OS}(2)$ is not
a notational artifact of how $2\times2$ matrices happen to be written; it
is a true, dimension-specific fact that is simply false one dimension
higher. No change of representation at $d=2$ would have revealed this —
the only diagnostic is varying the parameter ($d$) and recomputing, which
is a genuinely different check than asking whether a result survives
recoordinatization.

Both hazards share a single moral — a pattern that holds in a special case
can masquerade as the general rule if the special case is never
stress-tested against the parameter, or representation, that made it
special — but the stress test differs by hazard, and treating them as
interchangeable would leave a reader checking only one of the two ways a
result can be an artifact rather than a fact.
\end{remark}
